\documentclass[]{beamer}

\usetheme{metropolis}
\metroset{progressbar=frametitle}
\usefonttheme[onlymath]{serif}

\usepackage[none]{hyphenat}
\usepackage{hyperref}
\usepackage{amsmath}
\usepackage{amsfonts}
\usepackage{amssymb}
\usepackage[normalem]{ulem}
\usepackage{graphicx}
\usepackage{tikz}
\usepackage{minted}
\subtitle{Recommendation Systems}
\author{Maxim Borisyak}

\institute{Constructor University Bremen}

\usepackage{algorithm}
\usepackage{booktabs}
\usepackage{algpseudocode}
\usepackage{setspace}
\usepackage{framed}

\DeclareMathOperator*{\E}{\mathbb{E}}

\DeclareMathOperator*{\var}{\mathbb{D}}
\newcommand\D[1]{\var\left[ #1 \right]}

\newcommand\dmid{\,\|\,}

\DeclareMathOperator*{\argmin}{\mathrm{arg\,min}}
\DeclareMathOperator*{\argmax}{\mathrm{arg\,max}}




\begin{document}



\title{Machine Learning and Data Mining}

\begin{frame}[plain]
\titlepage
\end{frame}

\section{Recommendation Systems}

\begin{frame}[fragile]
\frametitle{What are recommendation systems?}
\textbf{Goal}: predict user preferences for items they have not yet interacted with.

\textbf{Applications}:
\begin{itemize}
  \item movies: Netflix, YouTube;
  \item music: Spotify, Last.fm;
  \item products: Amazon, Alibaba;
  \item news, social media, advertisements.
\end{itemize}

\textbf{Key challenge}: the rating matrix is extremely sparse --- most users interact with a tiny fraction of all items.

\end{frame}

\begin{frame}[fragile]
\frametitle{Problem formulation}
\textbf{Setting}:
\begin{itemize}
  \item $m$ users: $u \in \{1, \dots, m\}$;
  \item $n$ items: $i \in \{1, \dots, n\}$;
  \item rating matrix $R \in \mathbb{R}^{m \times n}$, where $R_{ui}$ is the rating of user $u$ for item $i$;
  \item $\Omega \subseteq \{1,\dots,m\} \times \{1,\dots,n\}$ --- set of observed (user, item) pairs.
\end{itemize}

\textbf{Goal}: estimate $R_{ui}$ for all $(u, i) \notin \Omega$.

\end{frame}

\begin{frame}[fragile]
\frametitle{Rating matrix}
\begin{equation*}
R = \begin{pmatrix}
  5 & \cdot & 3 & \cdot & 1 \\
  4 & \cdot & \cdot & 1 & \cdot \\
  \cdot & 2 & \cdot & 4 & \cdot \\
  \cdot & \cdot & 5 & \cdot & 2 \\
  1 & \cdot & 4 & \cdot & 5
\end{pmatrix}
\end{equation*}

\begin{itemize}
  \item rows = users; columns = items;
  \item $R_{ui} \in \{1,\dots,5\}$ or missing ($\cdot$);
  \item goal: fill in the missing entries.
\end{itemize}

\end{frame}

\begin{frame}[fragile]
\frametitle{Types of feedback}
\textbf{Explicit feedback}:
\begin{itemize}
  \item user actively rates an item (1--5 stars, thumbs up/down);
  \item sparse but directly expresses preference;
  \item examples: movie ratings, book reviews.
\end{itemize}

\textbf{Implicit feedback}:
\begin{itemize}
  \item inferred from user behavior;
  \item denser but noisier --- absence $\neq$ dislike;
  \item examples: clicks, purchases, watch time, skips.
\end{itemize}

\end{frame}

\begin{frame}[fragile]
\frametitle{Approaches to recommendation}
\textbf{Content-based filtering}:
\begin{itemize}
  \item recommend items similar to those the user liked;
  \item requires item features (genre, author, keywords);
  \item no interaction with other users.
\end{itemize}

\textbf{Collaborative filtering}:
\begin{itemize}
  \item exploit patterns across all users and items;
  \item no item features required;
  \item "users who liked $X$ also liked $Y$".
\end{itemize}

\textbf{Hybrid methods}: combine content-based and collaborative filtering.

\end{frame}

\section{Collaborative Filtering}

\begin{frame}[fragile]
\frametitle{Key idea}
\begin{center}
\textit{``Tell me who your friends are, and I will tell you what you like.''}
\end{center}

\vspace{4mm}

\textbf{Assumptions}:
\begin{itemize}
  \item users with similar past behavior will have similar future preferences;
  \item items that appealed to similar users will appeal to the target user.
\end{itemize}

\textbf{Two families}:
\begin{itemize}
  \item \textbf{user-based}: find users similar to the target user, aggregate their ratings;
  \item \textbf{item-based}: find items similar to the target item, aggregate ratings of those items.
\end{itemize}

\end{frame}

\begin{frame}[fragile]
\frametitle{User-based collaborative filtering}
\textbf{Algorithm}:
\begin{enumerate}
  \item compute similarity $\mathrm{sim}(u, v)$ between target user $u$ and every other user $v$;
  \item select $k$ nearest neighbors $\mathcal{N}(u)$;
  \item predict the rating as a similarity-weighted average.
\end{enumerate}

\begin{equation*}
\hat{R}_{ui} = \bar{r}_u + \frac{\displaystyle\sum_{v \in \mathcal{N}(u)} \mathrm{sim}(u, v)\,(R_{vi} - \bar{r}_v)}{\displaystyle\sum_{v \in \mathcal{N}(u)} |\mathrm{sim}(u, v)|}
\end{equation*}

where $\bar{r}_u = \frac{1}{|\Omega_u|} \sum_{j \in \Omega_u} R_{uj}$ is the mean rating of user $u$.

\end{frame}

\begin{frame}[fragile]
\frametitle{Similarity measures}
\textbf{Cosine similarity} (treats missing as zero):
\begin{equation*}
\mathrm{sim}_{\cos}(u, v) = \frac{\mathbf{r}_u \cdot \mathbf{r}_v}{\|\mathbf{r}_u\|\,\|\mathbf{r}_v\|}
\end{equation*}

\textbf{Pearson correlation} (uses co-rated items only):
\begin{equation*}
\mathrm{sim}_{\mathrm{P}}(u, v) =
  \frac{\displaystyle\sum_{i \in \Omega_{uv}} (R_{ui} - \bar{r}_u)(R_{vi} - \bar{r}_v)}
       {\sqrt{\displaystyle\sum_{i \in \Omega_{uv}} (R_{ui} - \bar{r}_u)^2}\;
        \sqrt{\displaystyle\sum_{i \in \Omega_{uv}} (R_{vi} - \bar{r}_v)^2}}
\end{equation*}

where $\Omega_{uv}$ = items rated by both $u$ and $v$.

\end{frame}

\begin{frame}[fragile]
\frametitle{Item-based collaborative filtering}
\textbf{Algorithm}:
\begin{enumerate}
  \item compute similarity $\mathrm{sim}(i, j)$ between target item $i$ and every other item $j$;
  \item select $k$ nearest neighbor items $\mathcal{N}(i)$;
  \item predict the rating using the user's ratings of neighboring items.
\end{enumerate}

\begin{equation*}
\hat{R}_{ui} = \frac{\displaystyle\sum_{j \in \mathcal{N}(i)} \mathrm{sim}(i, j)\,R_{uj}}
                    {\displaystyle\sum_{j \in \mathcal{N}(i)} |\mathrm{sim}(i, j)|}
\end{equation*}

\textbf{Key difference}: similarities are between items, not users.

\end{frame}

\begin{frame}[fragile]
\frametitle{User-based vs item-based}
\begin{center}
\begin{tabular}{lll}
\toprule
\textbf{Property} & \textbf{User-based} & \textbf{Item-based} \\
\midrule
Similarity space  & users ($m \times m$)   & items ($n \times n$) \\
Stability         & low (user tastes change) & high (catalog is stable) \\
Sparsity handling & struggles              & more robust \\
Explanation       & ``users like you liked\dots'' & ``similar to item $X$'' \\
\bottomrule
\end{tabular}
\end{center}

\vspace{3mm}
\textbf{In practice}: item-based CF is often preferred --- item similarities are more stable and can be precomputed offline.

\end{frame}

\begin{frame}[fragile]
\frametitle{Limitations of neighborhood methods}
\textbf{Sparsity}:
\begin{itemize}
  \item most users rate very few items;
  \item little overlap $\Omega_{uv}$ between user pairs.
\end{itemize}

\textbf{Scalability}:
\begin{itemize}
  \item all-pairs similarity: $O(m^2)$ or $O(n^2)$;
  \item prohibitive for millions of users/items.
\end{itemize}

\textbf{Cold start}:
\begin{itemize}
  \item new users or items have no ratings;
  \item no similarity can be computed.
\end{itemize}

\textbf{Quality ceiling}:
\begin{itemize}
  \item neighborhood heuristics are not globally consistent;
  \item matrix factorization methods typically outperform them.
\end{itemize}

\end{frame}

\section{Matrix Factorization}

\begin{frame}[fragile]
\frametitle{Latent factor models}
\textbf{Intuition}: ratings are driven by a small number of hidden factors.

\begin{itemize}
  \item each user $u$ is described by a latent vector $\mathbf{p}_u \in \mathbb{R}^k$;
  \item each item $i$ is described by a latent vector $\mathbf{q}_i \in \mathbb{R}^k$;
  \item predicted rating:
\begin{equation*}
\hat{R}_{ui} = \mathbf{p}_u^\top \mathbf{q}_i
\end{equation*}

\end{itemize}
\textbf{Examples of latent factors} (movies):
\begin{itemize}
  \item seriousness vs.\ humor;
  \item action-oriented vs.\ dialogue-driven;
  \item mainstream vs.\ arthouse.
\end{itemize}

\end{frame}

\begin{frame}[fragile]
\frametitle{Matrix factorization}
Approximate the full rating matrix as a product of two low-rank matrices:

\begin{equation*}
R \approx P Q^\top, \qquad P \in \mathbb{R}^{m \times k},\quad Q \in \mathbb{R}^{n \times k}
\end{equation*}

where $k \ll \min(m, n)$ is the number of latent factors.


\vfill

\begin{center}
\begin{tabular}{ccccc}
$R$ & $\approx$ & $P$ & $\cdot$ & $Q^\top$ \\[2mm]
$m \times n$ & & $m \times k$ & & $k \times n$ \\
\end{tabular}
\end{center}

\end{frame}

\begin{frame}[fragile]
\frametitle{Loss function}
\textbf{Objective} (minimize reconstruction error on observed entries only):

\begin{equation*}
\min_{P,\,Q}\; \sum_{(u,i)\in\Omega} \left(R_{ui} - \mathbf{p}_u^\top \mathbf{q}_i\right)^2
  + \lambda\!\left(\|P\|_F^2 + \|Q\|_F^2\right)
\end{equation*}

\begin{itemize}
  \item sum restricted to \textbf{observed} entries in $\Omega$;
  \item $\ell_2$ regularization with strength $\lambda > 0$ prevents overfitting;
  \item \textbf{non-convex} jointly in $(P, Q)$.
\end{itemize}

\end{frame}

\begin{frame}[fragile]
\frametitle{Bias terms}
\begin{equation*}
\hat{R}_{ui} = \mu + b_u + b_i + \mathbf{p}_u^\top \mathbf{q}_i
\end{equation*}

where:
\begin{itemize}
  \item $\mu$ --- global mean rating;
  \item $b_u$ --- user bias (some users rate higher overall);
  \item $b_i$ --- item bias (some items are rated higher overall).
\end{itemize}

Note: $\mathbf{p}'_u = (1, b_u, \mathbf{p}_u)^\top$, $\mathbf{q}'_i = (b_i, 1, \mathbf{q}_i)^\top$ --- a special case of the dot product model.

\end{frame}

\begin{frame}[fragile]
\frametitle{Why not just use SVD?}
\textbf{Truncated SVD} minimizes:
\begin{equation*}
\min_{\mathrm{rank}\,k}\; \|R - \hat{R}\|_F^2
\end{equation*}

over \textbf{all} entries --- but $R$ has missing values!

\textbf{Naive workaround}: impute missing entries with 0 or $\bar{r}$.

\textbf{Problems}:
\begin{itemize}
  \item imputing zeros biases toward low ratings;
  \item imputing the mean ignores structure;
  \item the resulting factorization fits the imputed values, not the true ones.
\end{itemize}

\textbf{Correct approach}: optimize only over observed entries $\Omega$.

\end{frame}

\section{Alternating Least Squares}

\begin{frame}[fragile]
\frametitle{ALS: motivation}
The combined loss is non-convex in $(P, Q)$:

\begin{equation*}
\mathcal{L}(P, Q) = \sum_{(u,i)\in\Omega}(R_{ui} - \mathbf{p}_u^\top \mathbf{q}_i)^2
  + \lambda\!\left(\|P\|_F^2 + \|Q\|_F^2\right)
\end{equation*}

But it is \textbf{convex} in $P$ alone when $Q$ is fixed, and vice versa.

\textbf{ALS strategy}: alternate between closed-form updates of $P$ and $Q$.

\end{frame}

\begin{frame}[fragile]
\frametitle{ALS: update for user factors}
Fix $Q$, minimize over each $\mathbf{p}_u$ independently.

Let $\Omega_u = \{i : (u,i) \in \Omega\}$ be the items rated by user $u$:

\begin{equation*}
\mathbf{p}_u = \left(Q_u^\top Q_u + \lambda I\right)^{-1} Q_u^\top \mathbf{r}_u
\end{equation*}

where:
\begin{itemize}
  \item $Q_u \in \mathbb{R}^{|\Omega_u| \times k}$ --- rows of $Q$ for items in $\Omega_u$;
  \item $\mathbf{r}_u \in \mathbb{R}^{|\Omega_u|}$ --- observed ratings of user $u$.
\end{itemize}

This is a \textbf{ridge regression} problem per user!

\end{frame}

\begin{frame}[fragile]
\frametitle{ALS: update for item factors}
Fix $P$, minimize over each $\mathbf{q}_i$ independently.

Let $\Omega_i = \{u : (u,i) \in \Omega\}$ be the users who rated item $i$:

\begin{equation*}
\mathbf{q}_i = \left(P_i^\top P_i + \lambda I\right)^{-1} P_i^\top \mathbf{r}_i
\end{equation*}

where:
\begin{itemize}
  \item $P_i \in \mathbb{R}^{|\Omega_i| \times k}$ --- rows of $P$ for users in $\Omega_i$;
  \item $\mathbf{r}_i \in \mathbb{R}^{|\Omega_i|}$ --- observed ratings for item $i$.
\end{itemize}

Symmetric to the user update --- also ridge regression!

\end{frame}

\begin{frame}[fragile]
\frametitle{ALS algorithm}
\begin{framed}
\begin{spacing}{1.5}
\begin{algorithmic}[1]
  \State Initialize $P$, $Q$ randomly (e.g.\ $\sim \mathcal{N}(0,\,\sigma^2)$)
  \Repeat
    \For{each user $u = 1,\dots,m$}
      \State $\mathbf{p}_u \leftarrow \left(Q_u^\top Q_u + \lambda I\right)^{-1} Q_u^\top \mathbf{r}_u$
    \EndFor
    \For{each item $i = 1,\dots,n$}
      \State $\mathbf{q}_i \leftarrow \left(P_i^\top P_i + \lambda I\right)^{-1} P_i^\top \mathbf{r}_i$
    \EndFor
  \Until{convergence}
\end{algorithmic}
\end{spacing}
\end{framed}

\end{frame}

\begin{frame}[fragile]
\frametitle{ALS: Parallelism:}
\begin{itemize}
  \item all user updates are independent $\Rightarrow$ trivially parallelizable;
  \item all item updates are independent $\Rightarrow$ trivially parallelizable;
  \item widely used in distributed settings (Apache Spark MLlib).
\end{itemize}

\textbf{Complexity per iteration}:
\begin{equation*}
O\!\left(|\Omega|\,k + (m + n)\,k^3\right)
\end{equation*}

\end{frame}

\begin{frame}[fragile]
\frametitle{Weighted ALS for implicit feedback}
\textbf{Implicit feedback} (Hu, Koren, Volinsky 2008):
\begin{itemize}
  \item let $R_{ui} = 1$ if user $u$ interacted with item $i$, else $0$;
  \item absence of interaction $\neq$ dislike --- only lower confidence.
\end{itemize}

\textbf{Weighted objective} (sum over \textbf{all} pairs):
\begin{equation*}
\min_{P,Q}\;\sum_{u=1}^{m}\sum_{i=1}^{n} c_{ui}\,(R_{ui} - \mathbf{p}_u^\top \mathbf{q}_i)^2
  + \lambda\!\left(\|P\|_F^2 + \|Q\|_F^2\right)
\end{equation*}

\textbf{Confidence}:
\begin{equation*}
c_{ui} = 1 + \alpha\;\mathrm{count}_{ui}
\end{equation*}
where $\mathrm{count}_{ui}$ is the number of interactions (plays, clicks, \dots).

\end{frame}

\begin{frame}[fragile]
\frametitle{Poisson factorization for implicit feedback}
\textbf{Model} (Gopalan, Hofman, Blei 2015):
\begin{equation*}
\mathrm{count}_{ui} \sim \mathrm{Poisson}\!\left(\mathbf{p}_u^\top \mathbf{q}_i\right), \quad \mathbf{p}_u, \mathbf{q}_i \geq 0
\end{equation*}

Negative Log-Likelihood:
\begin{equation*}
L(P, Q) = -\sum_{u,i} \left[ \mathbf{p}_u^\top \mathbf{q}_i - \mathrm{count}_{ui} \log(\mathbf{p}_u^\top \mathbf{q}_i) \right]
\end{equation*}

\begin{itemize}
  \item models \textbf{count data} (plays, clicks, views);
  \item Missing entries contribute via the $-\mathbf{p}_u^\top \mathbf{q}_i$ term;
\end{itemize}

\end{frame}

\begin{frame}[fragile]
\frametitle{Poisson factorization: log-linear parameterization}
Instead of constraining $\mathbf{p}_u, \mathbf{q}_i \geq 0$, parameterize the rate directly:
\begin{equation*}
\lambda_{ui} = \exp\!\left(\mathbf{p}_u^\top \mathbf{q}_i\right), \quad \mathbf{p}_u, \mathbf{q}_i \in \mathbb{R}^k
\end{equation*}

The log-likelihood becomes:
\begin{equation*}
\sum_{u,i} \left[ \mathrm{count}_{ui} \cdot \mathbf{p}_u^\top \mathbf{q}_i - e^{\mathbf{p}_u^\top \mathbf{q}_i} \right]
\end{equation*}

\begin{itemize}
  \item $\mathbf{p}_u, \mathbf{q}_i$ are unconstrained --- standard SGD applies directly;
  \item Gradients: $\nabla_{\mathbf{p}_u} = \sum_i (\mathrm{count}_{ui} - \lambda_{ui})\,\mathbf{q}_i$,
  \item ... symmetrically for $\mathbf{q}_i$.
\end{itemize}

\end{frame}

\begin{frame}[fragile]
\frametitle{Evaluation: explicit feedback}
\textbf{Root Mean Squared Error}:
\begin{equation*}
\mathrm{RMSE} = \sqrt{\frac{1}{|\Omega_{\mathrm{test}}|}
  \sum_{(u,i)\in\Omega_{\mathrm{test}}} \left(R_{ui} - \hat{R}_{ui}\right)^2}
\end{equation*}

\textbf{Mean Absolute Error}:
\begin{equation*}
\mathrm{MAE} = \frac{1}{|\Omega_{\mathrm{test}}|}
  \sum_{(u,i)\in\Omega_{\mathrm{test}}} |R_{ui} - \hat{R}_{ui}|
\end{equation*}

\end{frame}

\begin{frame}[fragile]
\frametitle{Evaluation: ranking metrics}
For implicit feedback or top-$N$ recommendation:

\textbf{Precision\@$N$}:
\begin{equation*}
\mathrm{Precision@}N = \frac{|\text{relevant items in top-}N|}{N}
\end{equation*}

\textbf{NDCG\@$N$} (Normalized Discounted Cumulative Gain):
\begin{equation*}
\mathrm{DCG@}N = \sum_{j=1}^{N} \frac{\mathrm{rel}_j}{\log_2(j+1)}, \qquad
\mathrm{NDCG@}N = \frac{\mathrm{DCG@}N}{\mathrm{IDCG@}N}
\end{equation*}

where $\mathrm{IDCG}$ is the ideal (perfect) ranking.

\end{frame}

\begin{frame}[fragile]
\frametitle{Practical considerations}
\textbf{Choosing rank $k$}:
\begin{itemize}
  \item too small: underfits, misses structure;
  \item too large: overfits, slow to train;
  \item typical range: $k \in \{10, 50, 200\}$.
\end{itemize}

\textbf{Choosing $\lambda$}:
\begin{itemize}
  \item cross-validation on held-out ratings;
  \item typical range: $\lambda \in \{0.01, 0.1, 1.0\}$.
\end{itemize}

\textbf{Initialization}:
\begin{itemize}
  \item small random values: $P, Q \sim \mathcal{N}(0, k^{-1/2})$;
  \item avoids large initial predictions.
\end{itemize}

\end{frame}

\begin{frame}[fragile]
\frametitle{Cold start problem}
\textbf{New user} (no interaction history):
\begin{itemize}
  \item cannot compute $\mathbf{p}_u$ directly;
  \item solutions: ask for initial ratings (onboarding), use demographic features, or start from the global mean vector.
\end{itemize}

\textbf{New item} (no ratings yet):
\begin{itemize}
  \item cannot compute $\mathbf{q}_i$ directly;
  \item solutions: initialize from item content features, use a separate content-based model.
\end{itemize}

\textbf{Fundamental limitation} of pure collaborative filtering.

\end{frame}

\begin{frame}[fragile]
\frametitle{Summary}
Recommendation systems:
\begin{itemize}
  \item Collaborative Filtering:
\begin{itemize}
  \item item-based;
  \item user-based;
\end{itemize}

  \item Matrix Factorization;
  \item Alternating Least Squares.
\end{itemize}

\vspace{3mm}
ALS-based matrix factorization dominates large-scale production systems (Netflix Prize winner, Spotify, etc.).
\end{frame}




\end{document}
